\section{The \texttt{KN\_xrl\_process} McXtrace Component}
A component implementing the Klein-Nishina cross section as a Union physics process

\subsection*{Identification}
\begin{itemize}
  \item \textbf{Author:} Mads Bertelsen and Erik B Knudsen
  \item \textbf{Origin:} ESS DMSC \& DTU Physics \& United Neux
  \item \textbf{Date:} 20.08.15
\end{itemize}

\subsection*{Description}
This Union\_process is based on the Incoherent.comp component originally written by Kim Lefmann and Kristian Nielsen

Part of the Union components, a set of components that work together and thus separates geometry and physics within McXtrace. The use of this component requires other components to be used.

1) One specifies a number of processes using process components like this one 2) These are gathered into material definitions using Union\_make\_material 3) Geometries are placed using Union\_box / Union\_cylinder, assigned a material 4) A Union\_master component placed after all of the above

Only in step 4 will any simulation happen, and per default all geometries defined before the master, but after the previous will be simulated here.

There is a dedicated manual available for the Union\_components

Algorithm: Described elsewhere

\subsection*{Input parameters}
Parameters in \textbf{boldface} are required; the others are optional.

\begin{longtable}{p{0.22\textwidth}p{0.12\textwidth}p{0.46\textwidth}p{0.14\textwidth}}
\toprule
\textbf{Name} & \textbf{Unit} & \textbf{Description} & \textbf{Default} \\
\midrule
\endhead
density & g/cm$^{3}$ & Nominal density of material. & 0 \\
init & str & Name of Union inititaliser component & "init" \\
\bottomrule
\end{longtable}

\subsection*{Links}
\begin{itemize}
  \item Component source code found in file \texttt{KN\_xrl\_process.comp}.
  \item The test/example instrument \htmladdnormallink{Test\_Phonon.instr}{../examples/Test\_Phonon.instr}.
\end{itemize}
\IfFileExists{union/KN_xrl_process_static.tex}{\input{union/KN_xrl_process_static.tex}}{}